\subsection{What Is Pragmatics}
\label{sec:what-is-pragmatics}
\slidesrc{Pragmatics: 1--3}

\begin{nbdef}{Pragmatics}
\textbf{Pragmatics} is a subfield of linguistics and semiotics that studies the ways in
which \textbf{context contributes to meaning}. It encompasses speech act theory,
conversational implicature, talk in interaction and other approaches to language
behaviour in philosophy, sociology, linguistics and anthropology.

Pragmatics studies how the transmission of meaning depends not only on the structural
and linguistic knowledge (grammar, lexicon, \ldots) of speaker and listener, but also on
\begin{itemize}[nosep]
  \item the \textbf{context of the utterance};
  \item any \textbf{pre-existing knowledge} about those involved;
  \item the \textbf{inferred intent} of the speaker;
  \item and other factors.
\end{itemize}
\end{nbdef}

\textbf{Semiotics} is the study of meaning-making: the study of sign processes
(\emph{semiosis}) and of meaningful communication. Words are one kind of sign; gestures,
pictures, traffic lights are others.

\begin{nbexample}{same words, different meanings}
``Can you pass the salt?'' Syntactically it is a question; semantically it asks about an
ability (are you able to pass the salt?). Pragmatically, at a dinner table, it is a
\emph{request}: answering ``yes'' without moving would be absurd. ``It's cold in
here'' can be a statement, a complaint, or a request to close the window, depending on
who says it to whom, and where. Only the context tells.
\end{nbexample}

\begin{nbwhy}
A chatbot that understands every word but not the intention behind it answers ``Yes, I
can'' to ``Can you tell me the time?''. Recognising intentions (\emph{intent detection})
is a central task of dialogue systems, and the reason why the agent communication
languages of Section~\ref{sec:chatbdi} label every message with the act it performs.
\end{nbwhy}
